\section{The coupling calculation and the remaining existence argument}
\label{fin-next}

The finite local induction supplies actual admissible interaction
families at every stage. Its exact density bookkeeping is
\[
 \gamma_{j+1}=\gamma_j-
      \beta_{j+1}\bigl(\mathcal F_{j,g_j}(E_j)\bigr),
 \qquad \gamma_j=g_j^{-2}.
\]
A trajectory of one physical density must obey this equation.
The bounded magnitude of its increments does not determine their
sign or tune an arbitrarily long ultraviolet sequence to a fixed
interacting endpoint.

The parity theorem isolates a concrete leading functional,
\[
 D_M=\Delta_M+\mathcal Q_{{\rm det},M}.
\]
Exact infinite classical transport and the quantum input bound
reduce the drift on the \(O(b^2)\) amplitude envelope to
\[
 \beta_{j+1}(D_{M_j})+
 O\!\left(\tau(T(g_j))+g_j^2T(g_j)^6+
                         \frac{g_jT(g_j)}b\right).
\]
Here \(b>0\) is fixed before sending the coupling to zero;
constants may depend on that fixed field radius and geometry.
This expression concerns the infinite quotient, with the
finite cutoff-extraction distinction retained above.
The leading Maxwell coefficient must be evaluated using the actual
stationary background, pivot coordinates, gauge normalization and
retained Gaussian. Every implicit background derivative belongs
to that calculation. A universal perturbative coefficient from
another scheme cannot substitute for this identification.

The next measure-theoretic steps remain substantial. The homogeneous
local germ must be represented on the required gauge-orbit domains,
and mixed and large-field histories must be incorporated with their
exact measure identities. Uniform source estimates must then produce
a nontrivial continuum observable hierarchy with the symmetry,
positivity and reconstruction properties specified in Notebook 01.
No conclusion about those steps follows merely from changing
interaction representatives or bounding a finite local integral.

\begin{thebibliography}{9}
\bibitem{priorfin}
\emph{Toward a four-dimensional Yang--Mills construction},
research Notebooks 15--30 in this workspace. The finite ordinary
minimum, common-domain local quantum map, total Ward extraction
and quantum quotient estimates are the principal inputs.
\bibitem{balaban109fin}
T.~Ba\l aban,
\emph{Renormalization group approach to lattice gauge field
theories. I. Generation of effective actions in a small field
approximation and a coupling constant renormalization in four
dimensions},
Commun. Math. Phys. \textbf{109} (1987), 249--301.
\href{https://doi.org/10.1007/BF01215223}{Publisher record}.
\bibitem{balaban116fin}
T.~Ba\l aban,
\emph{Renormalization group approach to lattice gauge field
theories. II. Cluster expansions},
Commun. Math. Phys. \textbf{116} (1988), 1--22.
\href{https://doi.org/10.1007/BF01239022}{Publisher record}.
The finite quotient and representative arguments here are given
separately; the published induction theorem is not being imported
with altered hypotheses.
\bibitem{sourcefin}
E.~Sussman and collaborators,
\href{https://github.com/ethanwsussman/Yang--Mills-LLM-experiment}
{Yang--Mills-LLM-experiment}, deposited raw accounts at commit
\texttt{e392ef45013a581664bab76b90867170b30fe10e}.
\end{thebibliography}
\end{document}
